\frfilename{mt215.tex}
\versiondate{13.11.13}
\copyrightdate{2000}

\def\chaptername{Taksonomi ruang ukur}
\def\sectionname{Ruang $\sigma$-hingga dan prinsip penghabisan}

\newsection{215}

Saya menyisipkan bagian singkat ini untuk membahas beberapa fakta berguna
yang mungkin luput jika ditempatkan dalam salah satu bagian yang lebih panjang
di bab ini. Hampir semua penerapan teori ukuran melibatkan ruang
$\sigma$-hingga, sampai-sampai banyak penulis hanya membahas sepintas ruang
lainnya. Saya sendiri lebih suka menandai pentingnya konsep-konsep semacam itu
dengan menyatakan secara tegas teorema mana yang hanya berlaku bagi kelas
ruang yang lebih terbatas tersebut. Namun, beberapa fakta tentang ruang
$\sigma$-hingga memang perlu dipahami sejak awal. Dalam 215B saya memberikan
daftar sifat yang mencirikan ruang $\sigma$-hingga. Sebagian di antaranya lebih
mudah dipahami melalui prinsip penghabisan (215A). Sekaligus, saya menyertakan
sebuah fakta mendasar mengenai ruang ukur
tanpa atom (215D).

\leader{215A}{Prinsip \dvrocolon{penghabisan}}\cmmnt{ Pernyataan berikut
merupakan contoh penggunaan salah satu metode terpenting dalam teori ukuran.

\medskip

\noindent}{\bf Lema} Misalkan $(X,\Sigma,\mu)$ sembarang ruang ukur dan
$\Cal E\subseteq\Sigma$ suatu keluarga tak-kosong sedemikian sehingga
$\sup_{n\in\Bbb N}\mu F_n$ berhingga untuk setiap barisan tak-menurun
$\sequencen{F_n}$ dalam $\Cal E$.

(a) Terdapat suatu barisan tak-menurun $\sequencen{F_n}$ dalam $\Cal E$
sedemikian sehingga, untuk setiap $E\in\Sigma$, {\it salah satu} dari dua
hal berikut berlaku: terdapat $n\in\Bbb N$ sedemikian sehingga $E\cup F_n$
tidak termuat dalam satu pun anggota $\Cal E$; {\it atau}, dengan menetapkan
$F=\bigcup_{n\in\Bbb N}F_n$,

\Centerline{$\lim_{n\to\infty}\mu(E\setminus F_n)
=\mu(E\setminus F)=0$.}

\noindent Khususnya, jika $E\in\Cal E$ dan $E\supseteq F$, maka
$E\setminus F$ terabaikan.

(b) Jika $\Cal E$ terarah ke atas, maka terdapat suatu barisan tak-menurun
$\sequencen{F_n}$ dalam $\Cal E$ sedemikian sehingga, dengan menetapkan
$F=\bigcup_{n\in\Bbb N}F_n$, berlaku $\mu F=\sup_{E\in\Cal E}\mu E$ dan
$E\setminus F$ terabaikan untuk setiap $E\in\Cal E$; dengan demikian $F$
merupakan supremum esensial dari $\Cal E$ dalam $\Sigma$\cmmnt{ dalam
pengertian 211G}.

(c) Jika gabungan setiap barisan tak-menurun dalam $\Cal E$ merupakan anggota
$\Cal E$, maka terdapat $F\in\Cal E$ sedemikian sehingga $E\setminus F$
terabaikan untuk setiap $E\in\Cal E$ dengan $F\subseteq E$.

\proof{{\bf (a)} Pilih $\sequencen{F_n}$, $\sequencen{\Cal E_n}$, dan
$\sequencen{u_n}$ secara induktif sebagai berikut. Ambil $F_0$ sebagai
sembarang anggota $\Cal E$. Setelah $F_n\in\Cal E$ dipilih, tetapkan
$\Cal E_n=\{E:F_n\subseteq E\in\Cal E\}$ dan
$u_n=\sup\{\mu E:E\in\Cal E_n\}$ dalam $[0,\infty]$, lalu pilih
$F_{n+1}\in\Cal E_n$ sedemikian sehingga
$\mu F_{n+1}\ge\min(n,u_n-2^{-n})$; lanjutkan proses ini.

Perhatikan bahwa konstruksi ini menghasilkan suatu barisan tak-menurun
$\sequencen{F_n}$ dalam $\Cal E$. Karena
$\Cal E_{n+1}\subseteq\Cal E_n$ untuk setiap $n$, barisan
$\sequencen{u_n}$ tak-menaik dan memiliki limit $u$ dalam $[0,\infty]$.
Karena $\min(n,u-2^{-n})\le\mu F_{n+1}\le u_n$ untuk setiap $n$, berlaku
$\lim_{n\to\infty}\mu F_n=u$. Hipotesis kita mengenai $\Cal E$ sekarang
menunjukkan bahwa $u$ berhingga.

Jika $E\in\Sigma$ sedemikian sehingga untuk setiap $n\in\Bbb N$ terdapat
$E_n\in\Cal E$ dengan $E\cup F_n\subseteq E_n$, maka
$E_n\in\Cal E_n$, sehingga

\Centerline{$\mu F_n\le\mu(E\cup F_n)\le\mu E_n\le u_n$}

\noindent untuk setiap $n$, dan
$\lim_{n\to\infty}\mu(E\cup F_n)=u$. Ini berarti bahwa

\Centerline{$\mu(E\setminus F)
\le\lim_{n\to\infty}\mu(E\setminus F_n)
=\lim_{n\to\infty}\mu(E\cup F_n)-\mu F_n
=0$,}

\noindent sebagaimana dinyatakan. Khususnya, kesimpulan ini berlaku jika
$E\in\Cal E$ dan $E\supseteq F$.

\medskip

{\bf (b)} Ambil $\sequencen{F_n}$ dari (a). Jika $E\in\Cal E$, maka
(karena $\Cal E$ terarah ke atas) $E\cup F_n$ termuat dalam suatu anggota
$\Cal E$ untuk setiap $n\in\Bbb N$; oleh sebab itu alternatif kedua dalam
(a) harus berlaku, dan $E\setminus F$ terabaikan. Akibatnya,

\Centerline{$\sup_{E\in\Cal E}\mu E
\le\mu F
=\lim_{n\to\infty}\mu F_n
\le\sup_{E\in\Cal E}\mu E$,}

\noindent sehingga $\mu F=\sup_{E\in\Cal E}\mu E$.

Jika $G$ sembarang himpunan terukur sedemikian sehingga $E\setminus G$
terabaikan untuk setiap $E\in\Cal E$, maka $F_n\setminus G$ terabaikan
untuk setiap $n$, sehingga $F\setminus G$ terabaikan; jadi $F$ merupakan
supremum esensial bagi $\Cal E$.

\medskip

{\bf (c)} Ambil kembali $\sequencen{F_n}$ dari (a), dan tetapkan
$F=\bigcup_{n\in\Bbb N}F_n$. Hipotesis kita sekarang menyatakan bahwa
$F\in\Cal E$, sehingga kedua sifat yang dinyatakan terpenuhi.
}%end of proof of 215A

\leader{215B}{}\cmmnt{ Ruang $\sigma$-hingga begitu penting sehingga
fakta-fakta berikut layak dinyatakan secara terperinci.

\medskip

\noindent}{\bf Proposisi} Misalkan $(X,\Sigma,\mu)$ suatu ruang ukur
semihingga. Tuliskan $\Cal N$ untuk keluarga himpunan $\mu$-terabaikan
dan $\Sigma^f$ untuk keluarga himpunan terukur berukuran berhingga.
Pernyataan-pernyataan berikut ekuivalen:

(i) $(X,\Sigma,\mu)$ bersifat $\sigma$-hingga;

(ii) setiap keluarga saling lepas dalam $\Sigma^f\setminus\Cal N$
terhitung;

(iii) setiap keluarga saling lepas dalam $\Sigma\setminus\Cal N$
terhitung;

(iv) untuk setiap $\Cal E\subseteq\Sigma$ terdapat himpunan terhitung
$\Cal E_0\subseteq\Cal E$ sedemikian sehingga
$E\setminus\bigcup\Cal E_0$ terabaikan untuk setiap $E\in\Cal E$;

(v) untuk setiap $\Cal E\subseteq\Sigma$ yang tak-kosong dan terarah ke
atas, terdapat suatu barisan tak-menurun $\sequencen{F_n}$ dalam
$\Cal E$ sedemikian sehingga
$E\setminus\bigcup_{n\in\Bbb N}F_n$ terabaikan untuk setiap
$E\in\Cal E$;

(vi) untuk setiap $\Cal E\subseteq\Sigma$ yang tak-kosong, terdapat suatu
barisan tak-menurun $\sequencen{F_n}$ dalam $\Cal E$ sedemikian sehingga
$E\setminus\bigcup_{n\in\Bbb N}F_n$ terabaikan untuk setiap $E\in\Cal E$
yang memenuhi $E\supseteq F_n$ untuk setiap $n\in\Bbb N$;

(vii) {\it $\mu X=0$}, {\it atau} terdapat suatu ukuran
peluang $\nu$ pada $X$ dengan domain dan himpunan-himpunan terabaikan yang
sama dengan $\mu$;

(viii) terdapat suatu fungsi terukur dan terintegralkan
$f:X\to\ocint{0,1}$;

(ix) $\mu X=0$, atau terdapat suatu fungsi terukur
$f:X\to\ooint{0,\infty}$ sedemikian sehingga $\int fd\mu=1$.

\proof{{\bf (i)$\Rightarrow$(vii) dan (viii)} Jika $\mu X=0$, maka (vii)
langsung berlaku dan kita dapat mengambil $f=\chi X$ dalam (viii). Jika
tidak, misalkan $\sequencen{E_n}$ suatu barisan saling lepas dalam
$\Sigma^f$ yang menutupi $X$. Mudah dilihat bahwa terdapat suatu barisan
$\sequencen{\alpha_n}$ bilangan real yang semuanya positif sedemikian sehingga
$\sum_{n=0}^{\infty}\alpha_n\mu E_n=1$. Tetapkan
$\nu E=\sum_{n=0}^{\infty}\alpha_n\mu(E\cap E_n)$ untuk $E\in\Sigma$;
maka $\nu$ merupakan ukuran peluang dengan domain $\Sigma$ dan
himpunan-himpunan terabaikan yang sama dengan $\mu$. Selain itu,
$f=\sum_{n=0}^{\infty}\min(1,\alpha_n)\chi E_n$ merupakan fungsi terukur
dan terintegralkan yang bernilai positif di setiap titik.

\medskip

{\bf (vii)$\Rightarrow$(vi) dan (v)} Andaikan (vii), dan misalkan
$\Cal E$ suatu keluarga tak-kosong yang terdiri atas himpunan terukur. Jika
$\mu X=0$, maka (vi) dan (v) jelas benar. Jika tidak, misalkan $\nu$
suatu ukuran peluang dengan domain $\Sigma$ dan himpunan-himpunan
terabaikan yang sama dengan $\mu$. Karena
$\sup_{E\in\Cal E}\nu E\le 1$ berhingga, kita dapat menerapkan 215Aa
untuk memperoleh suatu barisan tak-menurun $\sequencen{F_n}$ dalam
$\Cal E$ sedemikian sehingga
$E\setminus\bigcup_{n\in\Bbb N}F_n$ terabaikan untuk setiap
$E\in\Cal E$ yang
memuat $\bigcup_{n\in\Bbb N}F_n$; dan jika $\Cal E$ terarah ke atas,
$E\setminus\bigcup_{n\in\Bbb N}F_n$ terabaikan untuk setiap $E\in\Cal E$,
sebagaimana dalam 215Ab.

\medskip

{\bf (vi)$\Rightarrow$(iv)} Andaikan (vi), dan misalkan $\Cal E$
sembarang subhimpunan dari $\Sigma$. Tetapkan

\Centerline{$\Cal H=\{\bigcup\Cal E_0:\Cal E_0\subseteq\Cal E$ terhitung$\}$.}

\noindent Menurut (vi), terdapat suatu barisan $\sequencen{H_n}$ dalam
$\Cal H$ sedemikian sehingga
$H\setminus\bigcup_{n\in\Bbb N}H_n$ terabaikan untuk setiap $H\in\Cal H$
yang memenuhi $H\supseteq H_n$ untuk setiap $n\in\Bbb N$. Sekarang setiap $H_n$
dapat dinyatakan sebagai $\bigcup\Cal E'_n$, dengan
$\Cal E'_n\subseteq\Cal E$ terhitung; dengan menetapkan
$\Cal E_0=\bigcup_{n\in\Bbb N}\Cal E'_n$, himpunan $\Cal E_0$
terhitung. Jika $E\in\Cal E$, maka
$E\cup\bigcup_{n\in\Bbb N}H_n=\bigcup(\{E\}\cup\Cal E_0)$ termasuk
dalam $\Cal H$ dan memuat setiap $H_n$, sehingga
$E\setminus\bigcup\Cal E_0=E\setminus\bigcup_{n\in\Bbb N}H_n$
terabaikan. Jadi $\Cal E_0$ mempunyai sifat yang diperlukan, dan (iv)
benar.

\medskip

{\bf (iv)$\Rightarrow$(iii)} Andaikan (iv). Jika $\Cal E$ suatu keluarga
saling lepas dalam $\Sigma\setminus\Cal N$, ambil suatu
$\Cal E_0\subseteq\Cal E$ terhitung sedemikian sehingga
$E\setminus\bigcup\Cal E_0$ terabaikan untuk setiap $E\in\Cal E$.
Maka $E=E\setminus\bigcup\Cal E_0$ terabaikan untuk setiap
$E\in\Cal E\setminus\Cal E_0$; tetapi ini berarti
$\Cal E\setminus\Cal E_0$ kosong, sehingga $\Cal E=\Cal E_0$ terhitung.

\medskip

{\bf (iii)$\Rightarrow$(ii)} langsung.

\medskip

{\bf (ii)$\Rightarrow$(i)} Andaikan (ii). Misalkan $\frak P$ himpunan
semua keluarga saling lepas yang merupakan subhimpunan
$\Sigma^f\setminus\Cal N$, dengan urutan $\subseteq$. Maka $\frak P$
merupakan himpunan terurut sebagian yang tak-kosong (karena
$\emptyset\in \frak P$); dan setiap
$\frak Q\subseteq \frak P$ yang tak-kosong dan terurut total mempunyai
suatu batas atas dalam $\frak P$. \Prf\ Tetapkan
$\Cal E=\bigcup \frak Q$, yakni gabungan semua keluarga saling lepas yang
merupakan anggota $\frak Q$. Jika $E\in\Cal E$, maka $E\in\Cal C$ untuk
suatu $\Cal C\in \frak Q$, sehingga $E\in\Sigma^f\setminus\Cal N$.
Jika $E$, $F\in\Cal E$ dan $E\ne F$, maka terdapat
$\Cal C$, $\Cal D\in \frak Q$ sedemikian sehingga $E\in\Cal C$ dan
$F\in\Cal D$; karena $\frak Q$ terurut total, salah satu dari $\Cal C$,
$\Cal D$ memuat yang lain, dan dalam kedua kasus
$\Cal C\cup\Cal D$ merupakan anggota $\frak Q$ yang memuat $E$ dan $F$.
Karena setiap anggota $\frak Q$ merupakan keluarga himpunan saling lepas,
$E\cap F=\emptyset$. Karena $E$ dan $F$ sembarang, $\Cal E$ merupakan
keluarga himpunan saling lepas, sehingga merupakan anggota $\frak P$.
Tentu saja
$\Cal C\subseteq\Cal E$ untuk setiap $\Cal C\in \frak Q$, sehingga
$\Cal E$ merupakan batas atas bagi $\frak Q$ dalam $\frak P$.\ \Qed

Menurut Lema Zorn (2A1M), $\frak P$ mempunyai suatu elemen maksimal,
sebut $\Cal E$. Menurut (ii), $\Cal E$ harus terhitung, sehingga
$\bigcup\Cal E\in\Sigma$. Sekarang
$H=X\setminus\bigcup\Cal E$ terabaikan. \Prf\Quer\ Andaikan sebaliknya.
Karena $(X,\Sigma,\mu)$ semihingga, terdapat suatu himpunan $G$ berukuran
berhingga sedemikian sehingga $G\subseteq H$ dan $\mu G>0$, yakni
$G\in\Sigma^f\setminus\Cal N$ dan $G\cap E=\emptyset$ untuk setiap
$E\in\Cal E$. Namun ini berarti bahwa $\{G\}\cup\Cal E$ merupakan
anggota $\frak P$ yang lebih besar daripada $\Cal E$, padahal hal itu
mustahil.\ \Bang\Qed

Misalkan $\sequencen{X_n}$ suatu barisan yang mencakup semua anggota
$\Cal E\cup\{H\}$. Maka $\sequencen{X_n}$ merupakan penutup $X$ oleh
suatu barisan himpunan terukur berukuran berhingga, sehingga
$(X,\Sigma,\mu)$ bersifat $\sigma$-hingga.

\medskip

{\bf (v)$\Rightarrow$(i)} Jika (v) benar, maka terdapat suatu barisan
$\sequencen{E_n}$ dalam $\Sigma^f$ sedemikian sehingga
$E\setminus\bigcup_{n\in\Bbb N}E_n$ terabaikan untuk setiap
$E\in\Sigma^f$. Karena $\mu$ semihingga,
$X\setminus\bigcup_{n\in\Bbb N}E_n$ harus terabaikan; jadi $X$ ditutupi
oleh suatu keluarga terhitung yang terdiri atas himpunan berukuran berhingga
dan $\mu$ bersifat $\sigma$-hingga.

\medskip

{\bf (viii)$\Rightarrow$(ix)} Jika $\mu X=0$, pernyataan ini langsung.
Jika tidak, apabila $f$ suatu fungsi terukur dan terintegralkan yang bernilai
positif di setiap titik, maka $c=\int f>0$ (122Rc), sehingga $\Bover1cf$
merupakan fungsi terukur yang bernilai positif di setiap titik dan berintegral
$1$.

{\bf (ix)$\Rightarrow$(i)} Jika $f:X\to\ooint{0,\infty}$ terukur dan
terintegralkan, maka
$\sequencen{\{x:f(x)\ge 2^{-n}\}}$ merupakan barisan himpunan berukuran
berhingga yang menutupi $X$.
}%end of proof of 215B

\leader{215C}{Akibat} Misalkan $(X,\Sigma,\mu)$ suatu ruang ukur
$\sigma$-hingga, dan andaikan $\Cal E\subseteq\Sigma$ sembarang keluarga
tak-kosong.

(a) Terdapat suatu barisan tak-menurun $\sequencen{F_n}$ dalam $\Cal E$
sedemikian sehingga, untuk setiap $E\in\Sigma$, {\it salah satu} dari
dua hal berikut berlaku: terdapat $n\in\Bbb N$ sedemikian sehingga
$E\cup F_n$ tidak termuat dalam satu pun anggota $\Cal E$; {\it atau}
$E\setminus\bigcup_{n\in\Bbb N}F_n$ terabaikan.

(b) Jika $\Cal E$ terarah ke atas, maka terdapat suatu barisan tak-menurun
$\sequencen{F_n}$ dalam $\Cal E$ sedemikian sehingga
$\bigcup_{n\in\Bbb N}F_n$ merupakan supremum esensial dari $\Cal E$
dalam $\Sigma$.

(c) Jika gabungan setiap barisan tak-menurun dalam $\Cal E$ merupakan anggota
$\Cal E$, maka terdapat $F\in\Cal E$ sedemikian sehingga
$E\setminus F$ terabaikan untuk setiap $E\in\Cal E$ dengan $F\subseteq E$.

\proof{ Menurut 215B, terdapat suatu ukuran berhingga total $\nu$ pada
$X$ dengan himpunan-himpunan terukur dan himpunan-himpunan terabaikan yang
sama dengan $\mu$. Karena $\sup_{E\in\Cal E}\nu E$ berhingga, kita dapat
menerapkan 215A pada $\nu$ untuk memperoleh hasil-hasil tersebut.
}%end of proof of 215C

\leader{215D}{}\cmmnt{ Sebagai contoh lain penggunaan prinsip penghabisan,
saya memberikan sebuah fakta mendasar mengenai ruang ukur tanpa atom.

\medskip

\noindent}{\bf Proposisi} Misalkan $(X,\Sigma,\mu)$ suatu ruang ukur tanpa
atom. Jika $E\in\Sigma$ dan $0\le\alpha\le\mu E<\infty$, maka terdapat
$F\in\Sigma$ sedemikian sehingga $F\subseteq E$ dan $\mu F=\alpha$.

\proof{{\bf (a)} Kita perlu mengetahui bahwa jika $G\in\Sigma$
tak-terabaikan dan $n\in\Bbb N$, maka terdapat $H\subseteq G$ sedemikian
sehingga $0<\mu H\le 2^{-n}\mu G$. \Prf\ Kita buktikan dengan induksi
terhadap $n$. Untuk $n=0$ pernyataan ini langsung. Untuk langkah induksi
ke $n+1$, gunakan
hipotesis induksi untuk memperoleh $H\subseteq G$ sedemikian sehingga
$0<\mu H\le 2^{-n}\mu G$. Karena $\mu$ tanpa atom, terdapat
$H'\subseteq H$ sedemikian sehingga $\mu H'$ dan $\mu(H\setminus H')$
keduanya terdefinisi dan bukan nol. Setidaknya salah satu di antaranya
berukuran kurang dari atau sama dengan $\bover12\mu H$, sehingga
memberikan suatu subhimpunan dari $G$ yang berukuran positif dan tidak lebih
dari $2^{-n-1}\mu G$.\ \Qed

Akibatnya, jika $G\in\Sigma$ berukuran positif dan berhingga serta
$\epsilon>0$, maka terdapat suatu himpunan terukur $H\subseteq G$
sedemikian sehingga $0<\mu H\le\epsilon$.

\medskip

{\bf (b)} Misalkan $\Cal H$ adalah keluarga semua $H\in\Sigma$ sedemikian
sehingga $H\subseteq E$ dan $\mu H\le\alpha$. Jika $\sequencen{H_n}$
sembarang barisan tak-menurun dalam $\Cal H$, maka
$\mu(\bigcup_{n\in\Bbb N}H_n)=\lim_{n\to\infty}\mu H_n\le\alpha$,
sehingga $\bigcup_{n\in\Bbb N}H_n\in\Cal H$. Jadi 215Ac menunjukkan
bahwa terdapat $F\in\Cal H$ sedemikian sehingga $H\setminus F$
terabaikan untuk setiap $H\in\Cal H$ dengan $F\subseteq H$.
\Quer\ Andaikan sebaliknya bahwa $\mu F<\alpha$. Menurut (a),
terdapat $H\subseteq E\setminus F$ sedemikian sehingga
$0<\mu H\le\alpha-\mu F$. Namun dalam hal ini $H\cup F\in\Cal H$ dan
$\mu((H\cup F)\setminus F)>0$, suatu kemustahilan.\ \Bang

Dengan demikian kita telah menemukan himpunan $F$ yang sesuai.
}%end of proof of 215D

\leader{*215E}{}\cmmnt{ Salah satu sifat dasar ukuran Lebesgue ialah bahwa
himpunan bagian beranggota tunggal (dan karenanya himpunan bagian terhitung)
terabaikan.
Hal ini tentu berkaitan dengan kenyataan bahwa ukuran Lebesgue bersifat tanpa atom
(211Md). Mudah mengonstruksi ukuran yang membuat setiap himpunan beranggota
tunggal terabaikan tetapi tetap memiliki atom (misalnya ukuran
terhitung-koterhitung dalam 211R). Perlu sedikit siasat tambahan untuk
mengonstruksi ukuran tanpa atom yang di dalamnya tidak semua himpunan
beranggota tunggal terabaikan
(216Ye). Hasil berikut memberikan syarat-syarat yang mencegah hal tersebut.

\medskip

\noindent}{\bf Proposisi}\dvAnew{2015}
Misalkan $(X,\Sigma,\mu)$ suatu ruang ukur tanpa atom dan $x\in X$.

(a) Jika $\mu^*\{x\}$ berhingga, maka $\{x\}$ terabaikan.

(b) Jika $\mu$ memiliki himpunan-himpunan terabaikan yang ditentukan
secara lokal, maka $\{x\}$ terabaikan.

(c) Jika $\mu$ dapat dilokalkan, maka $\{x\}$ terabaikan.

\proof{{\bf (a)} \Quer\ Andaikan sebaliknya; maka
$0<\mu^*\{x\}<\infty$. Misalkan $E$ suatu himpunan yang memuat $x$
sedemikian sehingga $\mu E<2\mu^*\{x\}$. Menurut 215D, terdapat
$F\subseteq E$ sedemikian sehingga

\Centerline{$\mu F=\bover12\mu E=\mu(E\setminus F)<\mu^*\{x\}$.}

\noindent Jadi baik $F$ maupun $E\setminus F$ tidak mungkin memuat $x$;
padahal $x\in E$.\ \Bang

\medskip

{\bf (b)} Untuk setiap himpunan $E$ berukuran berhingga,
$E\cap\{x\}$ terabaikan; sebab, jika himpunan ini tidak kosong, maka
$x\in E$ dan kita dapat menerapkan (a). Karena $\mu$ memiliki
himpunan-himpunan terabaikan yang ditentukan secara lokal, $\{x\}$
terabaikan.

\medskip

{\bf (c)} Misalkan $\Cal E\subseteq\Sigma^f$ suatu keluarga maksimal
sedemikian sehingga $\mu E<\infty$ untuk setiap $E\in\Cal E$ dan
$\mu(E\cap F)=0$ setiap kali $E$, $F\in\Cal E$ berbeda. Untuk setiap
$E\in\Cal E$ dan $n\in\Bbb N$, terapkan 215D sebanyak $n$ kali untuk
memperoleh suatu partisi $E_{n0}$, $E_{n1},\ldots,E_{nn}$ dari $E$
sedemikian sehingga $\mu E_{ni}=\bover1{n+1}\mu E$ untuk setiap $i\le n$.
Selanjutnya, untuk $i\le n\in\Bbb N$, misalkan $H_{ni}$ suatu supremum
esensial dari $\{E_{ni}:E\in\Cal E\}$. Untuk setiap $n\in\Bbb N$ dan
$E\in\Cal E$,

\Centerline{$\mu(E\setminus\bigcup_{i\le n}H_{ni})
\le\sum_{i=0}^n\mu(E_{ni}\setminus H_{ni})=0$.}

\noindent Jadi, jika $F\in\Sigma$, $\mu F<\infty$, dan
$F\cap\bigcup_{i\le n}H_{ni}=\emptyset$, maka $\mu(E\cap F)=0$ untuk
setiap $E\in\Cal E$; karena $\Cal E$ maksimal, $F$ harus terabaikan.
Kita mengandaikan bahwa $\mu$ semihingga, sehingga
$X\setminus\bigcup_{i\le n}H_{ni}$ terabaikan.

Akibatnya, jika terdapat $n$ sedemikian sehingga
$x\notin\bigcup_{i\le n}H_{ni}$, maka $\{x\}$ terabaikan dan pembuktian
selesai. Tinjau
$H=\bigcap\{H_{ni}:i\le n\in\Bbb N$, $x\in H_{ni}\}$.
\Quer\ Jika $\mu H>0$, terdapat $F\subseteq H$ sedemikian sehingga
$0<\mu F<\infty$. Karena kemaksimalan $\Cal E$, terdapat
$E^*\in\Cal E$ sedemikian sehingga $\mu(F\cap E^*)>0$. Pilih $n$
sedemikian sehingga
$\mu(F\cap E^*)>\bover1{n+1}\mu E^*$. Terdapat $i\le n$ sedemikian
sehingga $x\in H_{ni}$; dengan demikian $F\subseteq H_{ni}$, sedangkan
$\mu(F\cap E^*_{ni})\le\bover1{n+1}\mu E^*$, sehingga
$F'=F\cap E^*\setminus E^*_{ni}$ berukuran positif. Tinjau
$G=H_{ni}\setminus F'$. Jika $E\in\Cal E$, maka salah satu dari dua hal
berikut berlaku: $E=E^*$ dan

\Centerline{$E_{ni}\setminus G=E^*_{ni}\setminus G
=E^*_{ni}\setminus H_{ni}$}

\noindent terabaikan; atau $E\ne E^*$ dan

\Centerline{$E_{ni}\setminus G
\subseteq(E_{ni}\setminus H_{ni})\cup(E\cap E^*)$}

\noindent terabaikan. Karena $H_{ni}$ merupakan supremum esensial dari
$\{E_{ni}:E\in\Cal E\}$, maka $H_{ni}\setminus G$ harus terabaikan dan
$F'$ terabaikan; ini mustahil.\ \Qed

Jadi $H$ merupakan himpunan terabaikan yang memuat $x$, dan dalam kasus
ini pun $\{x\}$ terabaikan.
}%end of proof of 215E

\exercises{\leader{215X}{Latihan dasar (a)}
%\spheader 215Xa
Misalkan $(X,\Sigma,\mu)$ suatu ruang ukur dan $\Phi$ suatu himpunan
tak-kosong yang terdiri atas fungsi-fungsi bernilai real yang terintegralkan
terhadap $\mu$, yang masing-masing memetakan $X$ ke $\Bbb R$. Andaikan
$\sup_{n\in\Bbb N}\int f_n$ berhingga
untuk setiap barisan $\sequencen{f_n}$ dalam $\Phi$ sedemikian sehingga
$f_n\leae f_{n+1}$ untuk setiap $n$. Tunjukkan bahwa terdapat suatu
barisan $\sequencen{f_n}$ dalam $\Phi$ sedemikian sehingga
$f_n\leae f_{n+1}$ untuk setiap $n$ dan, untuk setiap fungsi bernilai
real yang terintegralkan $f$ pada $X$, {\it salah satu}:
$f\leae\sup_{n\in\Bbb N}f_n$ {\it atau} terdapat $n\in\Bbb N$ sedemikian
sehingga tidak ada anggota $\Phi$ yang lebih besar atau sama
dengan $\max(f,f_n)$ hampir di mana-mana.
%215A

\sqheader 215Xb Misalkan $(X,\Sigma,\mu)$ suatu ruang ukur. (i) Andaikan
$\Cal E$ suatu subhimpunan $\Sigma$ yang tak-kosong dan terarah ke atas
sedemikian sehingga $c=\sup_{E\in\Cal E}\mu E$ berhingga. Tunjukkan
bahwa $E\setminus\bigcup_{n\in\Bbb N}F_n$ terabaikan untuk setiap
$E\in\Cal E$ dan setiap barisan $\sequencen{F_n}$ dalam $\Cal E$
sedemikian sehingga $\lim_{n\to\infty}\mu F_n=c$. (ii) Misalkan $\Phi$
suatu himpunan tak-kosong yang terdiri atas fungsi-fungsi terintegralkan
pada $X$
yang terarah ke atas dalam arti bahwa untuk setiap $f$, $g\in\Phi$
terdapat $h\in\Phi$ sedemikian sehingga $\max(f,g)\leae h$, serta andaikan
$c=\sup_{f\in\Phi}\int f$ berhingga. Tunjukkan bahwa
$f\leae\sup_{n\in\Bbb N}f_n$ untuk setiap $f\in\Phi$ dan setiap barisan
$\sequencen{f_n}$ dalam $\Phi$ sedemikian sehingga
$\lim_{n\to\infty}\int f_n=c$.
%215A

\spheader 215Xc Gunakan 215A untuk mempersingkat bukti 211Ld.

\spheader 215Xd Berikan contoh suatu ruang ukur (yang tidak semihingga)
$(X,\Sigma,\mu)$ yang memenuhi syarat (ii)-(iv) dalam 215B, tetapi tidak
memenuhi (i).
%215B

\sqheader 215Xe Misalkan $(X,\Sigma,\mu)$ suatu ruang ukur $\sigma$-hingga
tanpa atom. Tunjukkan bahwa untuk setiap $\epsilon>0$ terdapat suatu
barisan saling lepas $\sequencen{E_n}$ dari himpunan-himpunan terukur
yang masing-masing berukuran paling besar $\epsilon$ dan sedemikian
sehingga $X=\bigcup_{n\in\Bbb N}E_n$.
%215D

\spheader 215Xf Misalkan $(X,\Sigma,\mu)$ suatu ruang ukur tanpa atom
yang dapat dilokalkan secara ketat. Tunjukkan bahwa untuk setiap
$\epsilon>0$ terdapat suatu dekomposisi $\familyiI{X_i}$ dari $X$
sedemikian sehingga $\mu X_i\le\epsilon$ untuk setiap $i\in I$.

\leader{215Y}{Latihan lanjutan (a)}
%\spheader 215Ya
Misalkan $(X,\Sigma,\mu)$ suatu ruang ukur $\sigma$-hingga dan
$\langle f_{mn}\rangle_{m,n\in\Bbb N}$, $\sequence{m}{f_m}$, $f$
fungsi-fungsi terukur bernilai real yang didefinisikan hampir di
mana-mana dalam $X$ serta sedemikian sehingga
$\sequencen{f_{mn}}\to f_m$ hampir di mana-mana untuk setiap $m$ dan
$\sequence{m}{f_m}\to f$ hampir di mana-mana. Tunjukkan bahwa terdapat
suatu barisan naik ketat $\sequence{m}{n_m}$ dalam $\Bbb N$ sedemikian sehingga
$\sequence{m}{f_{m,n_m}}\to f$ hampir di mana-mana. (Bandingkan 134Yb.)
%215B

\spheader 215Yb Misalkan $(X,\Sigma,\mu)$ suatu ruang ukur
$\sigma$-hingga. Misalkan $\sequencen{f_n}$ suatu barisan fungsi terukur
bernilai real sedemikian sehingga $f=\lim_{n\to\infty}f_n$ didefinisikan
hampir di mana-mana dalam $X$. Tunjukkan bahwa terdapat suatu barisan
tak-menurun $\sequence{k}{X_k}$ yang terdiri atas subhimpunan terukur dari
$X$ sedemikian sehingga $\bigcup_{k\in\Bbb N}X_k$ koterabaikan dalam
$X$ dan $\sequencen{f_n}\to f$ seragam pada setiap $X_k$, dalam arti
bahwa untuk setiap $\epsilon>0$ terdapat $m\in\Bbb N$ sedemikian
sehingga $|f_j(x)-f(x)|$ terdefinisi dan kurang dari atau sama dengan
$\epsilon$ untuk setiap $j\ge m$ dan $x\in X_k$.

(Ini merupakan suatu versi Teorema Egorov.)
%215B

\spheader 215Yc Misalkan $(X,\Sigma,\mu)$ suatu ruang ukur berhingga
total dan $\sequencen{f_n}$, $f$ fungsi-fungsi terukur bernilai real
yang didefinisikan hampir di mana-mana dalam $X$. Tunjukkan bahwa
$\sequencen{f_n}\to f$ hampir di mana-mana jika dan hanya jika terdapat
suatu barisan $\sequencen{\epsilon_n}$ dari bilangan real yang semuanya
positif dan konvergen ke $0$, sedemikian sehingga

\Centerline{$\lim_{n\to\infty}\mu^*(\bigcup_{k\ge n}\{x:x\in\dom
f_k\cap\dom f,\,|f_k(x)-f(x)|\ge\epsilon_n\})=0$.}
%215B

\spheader 215Yd Temukan bukti langsung untuk
(v)$\Rightarrow$(vi) dalam 215B.
\Hint{untuk $\Cal E\subseteq\Sigma$, gunakan Lema Zorn untuk
memperoleh suatu $\Cal E'\subseteq\Cal E$ yang maksimal dan terurut total
sedemikian sehingga $E\symmdiff F\notin\Cal N$ untuk setiap
$E$, $F\in\Cal E'$ yang berbeda, lalu terapkan (v) pada $\Cal E'$.}

}%end of exercises

\endnotes{
\Notesheader{215}
Gagasan bersama yang mendasari 215A, 215B(vi), 215C, dan 215Xa
sebenarnya merupakan salah satu gagasan paling mendasar dalam teori
ukuran. Gagasan ini muncul dalam begitu banyak bentuk sehingga sering
kali lebih mudah membuktikan suatu penerapan langsung berdasarkan asas-asas dasar
daripada menjelaskan bagaimana penerapan tersebut dapat direduksi ke
versi-versi di sini. Namun, mulai sekarang saya akan berusaha menandai
penerapan semacam itu setiap kali muncul, dengan menyebut metode tersebut
(bukti 215Aa atau 215Xa) sebagai "prinsip penghabisan". Salah satu hal
yang mungkin patut dicatat di sini ialah konstruksi induktif barisan
$\sequencen{F_n}$ dalam bukti 215Aa. Setiap $F_{n+1}$ dipilih {\it
setelah} suku sebelumnya. Inilah yang memungkinkan kita, dalam bukti
215B(vii)$\Rightarrow$(vi), memperoleh langsung suatu barisan yang
sesuai $\sequencen{F_n}$. Dalam banyak penerapan (mulai dari penerapan yang
tentu paling penting dalam teori elementer, yakni Teorema Radon-Nikod\'ym
dalam \S232, ataupun bagian (i) dari bukti 211Ld),
penghalusan ini tidak diperlukan; kita berhadapan dengan suatu himpunan
yang terarah ke atas, seperti dalam 215B(v), dan dapat memilih seluruh
barisan $\sequencen{F_n}$ sekaligus, tanpa interaksi antarsuku, seperti
dalam 215Xb. Aksioma "pilihan bergantung", yang menyatakan bahwa kita
dapat mengonstruksi barisan suku demi suku, diketahui lebih kuat daripada
aksioma "pilihan terhitung", yang hanya menyatakan bahwa kita dapat
memilih objek dalam jumlah terhitung secara serentak.

Dalam 215B saya berusaha menunjukkan ciri-ciri utama sifat
$\sigma$-hingga; khususnya, sifat-sifat yang membedakan ukuran
$\sigma$-hingga dari ukuran lain yang dapat dilokalkan secara ketat.
Hasil ini, dalam arti tertentu, lebih abstrak daripada uraian lain dalam
bagian ini. Perhatikan bahwa hasil ini menggunakan
aksioma pilihan secara esensial dalam bentuk Lema Zorn. Dalam 134C saya
menulis satu paragraf untuk mengomentari perbedaan antara "pilihan
terhitung", yang diperlukan untuk segala sesuatu yang menyerupai teori
standar ukuran Lebesgue, dan aksioma pilihan penuh, yang relatif jarang
digunakan dalam teori elementer. Implikasi (ii)$\Rightarrow$(i) dalam
215B merupakan salah satu tempat kita benar-benar memerlukan
sesuatu yang melampaui pilihan terhitung. (Perlu kiranya saya catat bahwa
seluruh teori ruang ukur yang bukan $\sigma$-hingga tampak sangat ganjil
tanpa aksioma pilihan umum.) Perhatikan pula bahwa dalam 215B, bukti
(i)$\Rightarrow$(vii) dan (vii)$\Rightarrow$(vi) merupakan satu-satunya
bagian yang memunculkan sesuatu yang seremeh bilangan. Syarat (iii), (iv),
(v), dan (vi) saling terhubung dengan cara yang sama sekali tidak
berkaitan dengan teori ukuran, dan hanya melibatkan struktur
$(X,\Sigma,\Cal N)$. (Lihat 215Yd di sini, dan 316D-316E dalam Volume 3.)
Terdapat syarat-syarat serupa yang berkaitan dengan fungsi terukur,
bukan himpunan terukur; untuk suatu contoh yang cukup abstrak, lihat
241Ye.

Dalam 215Ya-215Yc terdapat tiga teorema standar lain mengenai barisan
yang konvergen hampir di mana-mana dan bergantung pada sifat
$\sigma$-hingga atau keberhinggaan total.
}%end of notes

\discrpage



